% BHU PYQ -- Auto-generated & error-corrected
\documentclass[11pt,a4paper]{article}

\usepackage[a4paper,top=2.5cm,bottom=2.5cm,left=2.8cm,right=2.8cm,headheight=14pt]{geometry}
\usepackage{amsmath,amssymb}
\usepackage[T1]{fontenc}
\usepackage[utf8]{inputenc}
\usepackage{lmodern}
\usepackage{microtype}
\usepackage{fancyhdr}
\usepackage{enumitem}
\usepackage{hyperref}
\usepackage{lastpage}
\usepackage{parskip}

\hypersetup{colorlinks=true,urlcolor=black,linkcolor=black,
  pdftitle={STB-301: Probability Theory and Statistical Methods -- BHU PYQ},pdfauthor={Banaras Hindu University}}

\pagestyle{fancy}\fancyhf{}
\renewcommand{\headrulewidth}{0.4pt}\renewcommand{\footrulewidth}{0.4pt}
\fancyhead[L]{\small   \textit{Banaras Hindu University}}
\fancyhead[R]{\small   \textit{STB-301: Probability Theory and Statistical }}
\fancyfoot[C]{\small Page  \thepage\ of \pageref{LastPage}}


\newlist{parts}{enumerate}{2}
\setlist[parts,1]{label=(\alph*),leftmargin=*,itemsep=5pt,topsep=3pt}
\setlist[parts,2]{label=(\roman*),leftmargin=*,itemsep=3pt,topsep=2pt}

\newcommand{\pts}[1]{\hfill   \textit{[\#1]}}


\begin{document}


\begin{center}
  {\large\bfseries BANARAS HINDU UNIVERSITY}\\[2pt]
  {\normalsize B.Sc. (Hons.) Semester III Examination, 2016-17}\\[6pt]
  \rule{0.8\linewidth}{0.5pt}\\[6pt]
  {\large\bfseries Paper No. STB-301: Probability Theory and Statistical Methods}\\[4pt]
  {\normalsize Time: 3 Hours\hfill Maximum Marks: 70}
\end{center}
\rule{\linewidth}{0.4pt}
\smallskip



\noindent   \textit{Instructions: Answer   \textbf{five} questions including Question~1 which is compulsory. Symbols have their usual meanings.}
\bigskip

\medskip


\noindent   \textbf{Question 1.}

\begin{parts}
  \item Explain the concept of sampling distribution with a suitable example. Suppose $X_1, X_2, \dots, X_n$ is a random sample of size $n$ from a population with finite mean $\mu$ and variance $\sigma^2$. Show that:
    $$E(ar{X}) = \mu \quad  \text{and} \quad V(ar{X}) = \frac{\sigma^2}{n}$$
    where $ar{X} = \frac{1}{n} \sum_{i=1}^n X_i$ is the sample mean.
  \item Let $X_1, X_2, \dots, X_n$ be a random sample from a Bernoulli distribution with parameter $p$. Show that the sampling distribution of $S = \sum_{i=1}^n X_i$ is a binomial distribution.
\end{parts}
\medskip\rule{\linewidth}{0.2pt}

\medskip


\noindent   \textbf{Question 2.}

\begin{parts}
  \item Define the probabilities of Type I and Type II errors. A sample of size 1 is taken from an exponential distribution with parameter $\lambda$. To test $H_0: \lambda = 1$ against $H_1: \lambda > 1$, suppose the critical region is considered as $x > 2$. Find the probabilities of Type I and Type II errors.
  \item Define the null and alternative hypotheses with examples. Explain with illustrations how to specify the null and alternative hypotheses in a hypothesis testing problem.
\end{parts}
\medskip\rule{\linewidth}{0.2pt}

\medskip


\noindent   \textbf{Question 3.}

\begin{parts}
  \item Write the density function of the $t$-distribution. Mention its distributional properties. Also, define Student's $t$-statistic and its use.
  \item What is $F$-distribution? Explain the testing procedure for testing the equality of two variances of two independent normal populations.
\end{parts}
\medskip\rule{\linewidth}{0.2pt}

\medskip


\noindent   \textbf{Question 4.}

\begin{parts}
  \item State the weak law of large numbers (WLLN). Let $X_1, X_2, \dots, X_n$ be a sequence of independent and identically distributed random variables with a common Bernoulli distribution with parameter $p$. Show that the sequence of sample means $\{ar{X}_n\}$ obeys the WLLN.
  \item State the central limit theorem (CLT). Suppose $X_1, X_2, \dots, X_n$ is a random sample of size $n$ from a population with finite mean $5$ and variance $16$. When $n$ is large, say $n = 40$, obtain the probability $P[ar{X} < 5]$.
\end{parts}
\medskip\rule{\linewidth}{0.2pt}

\medskip


\noindent   \textbf{Question 5.}

\begin{parts}
  \item Define order statistics with examples in real-life applications. Show that the $r$-th order statistic from a uniform distribution $U[0,1]$ follows a Beta distribution.
  \item Let $X_1, X_2, \dots, X_n$ be a random sample from a normal population with mean $\mu$ and variance $\sigma^2$. Show that the sampling distribution of the sample mean $ar{X}$ is also a normal distribution.
\end{parts}
\medskip\rule{\linewidth}{0.2pt}

\medskip


\noindent   \textbf{Question 6.}

\begin{parts}
  \item What are the advantages and disadvantages of non-parametric tests over parametric tests?
  \item Write the Mann-Whitney-Wilcoxon test statistic. To compare the mileage distributions of two brands of tires, the following mileages (in 1000 miles) were obtained for eight tires of each brand:
    
\begin{center}
      
\begin{tabular}{l|cccccccc}
        Brand A & 32.1 & 20.6 & 17.8 & 28.4 & 19.6 & 21.4 & 19.9 & 30.1 \\
        Brand B & 19.8 & 27.6 & 30.8 & 27.6 & 34.1 & 18.7 & 16.9 & 17.9
      \end{tabular}
    \end{center}
    Test the null hypothesis that the two samples come from the same population using the Mann-Whitney-Wilcoxon test.
\end{parts}
\medskip\rule{\linewidth}{0.2pt}

\medskip


\noindent   \textbf{Question 7.}
Write short notes on any three of the following:

\begin{parts}
  \item Uses of $\chi^2$ (chi-square) distribution.
  \item Median test.
  \item $p$-value.
  \item Test of independence of two attributes.
  \item Fisher's $Z$ transformation.
\end{parts}
\medskip\rule{\linewidth}{0.2pt}

\medskip


\noindent   \textbf{Question 8.}

\begin{parts}
  \item Explain the procedure of testing for the significance of the sample correlation coefficient $r$ in random sampling from a bivariate normal distribution.
  \item Nine students held their breath, once after breathing normally and relaxing for one minute, and once after hyperventilating for one minute. The table indicates how long (in seconds) they were able to hold their breath:
    
\begin{center}
      
\begin{tabular}{c|ccccccccc}
        Subject & A & B & C & D & E & F & G & H & I \\
        \hline
        Normal & 56.5 & 65.0 & 65.0 & 50.0 & 25.0 & 78.0 & 71.0 & 44.0 & 35.0 \\
        Hyperventilated & 87.0 & 91.0 & 85.0 & 91.0 & 75.0 & 28.0 & 122.0 & 66.0 & 38.0
      \end{tabular}
    \end{center}
    Mentioning the null and alternative hypotheses clearly, test the significance of the correlation coefficient between the two breath-holding times. (Given $t_{0.05, 7} = 2.365$ and $t_{0.05, 8} = 2.306$).
\end{parts}

\vfill
\rule{\linewidth}{0.4pt}

\begin{center}\small
  \href{https://bhu-exam.vercel.app/}{Click here to visit our website}\\[4pt]\href{https://me-aryan.vercel.app/}{Click here to visit our contributor}\\[4pt]\href{https://quashan-me.vercel.app/}{Click here to visit our contributor}
\end{center}

\end{document}